\chapter{Resultaten}
% TIP:
% Per deelvraag resultaten weergeven.
Dit hoofdstuk presenteert de resultaten van het onderzoek en geeft afzonderlijk antwoord op de vijf deelvragen.
De resultaten zijn gebaseerd op documentanalyse, analyse van de huidige systemen, literatuuronderzoek en gesprekken met medewerkers van Solvware.

\section{Resultaten voor Deelvraag 1}
% TIP: Analyse van huidige werkwijze.

\section{Resultaten voor Deelvraag 2}
\textbf{Welke functionele en niet-functionele eisen zijn noodzakelijk voor het ontwerp van een koppeling tussen SolvCRM en e-Boekhouden.nl?}

Op basis van documentanalyse, best practices voor API-integraties en gesprekken met de opdrachtgever zijn de noodzakelijke eisen vastgesteld. Deze eisen worden hieronder gepresenteerd als onderzoeksresultaten.

\subsection{Functionele eisen}

\subsubsection*{Klantinformatie}
Uit het onderzoek blijkt dat correcte synchronisatie van klantdata essentieel is voor een foutloze verwerking. Daarom zijn de volgende functionele eisen vastgesteld:

\begin{itemize}
    \item Automatische synchronisatie van klantinformatie tussen beide systemen.
    \item Mapping van klantgegevens tussen SolvCRM en e-Boekhouden.nl.
    \item Mogelijkheid om bestaande klanten te koppelen.
    \item Automatische aanmaak van nieuwe klanten in beide systemen.
    \item Verwijderde of gedeactiveerde klanten worden consistent doorgevoerd.
    \item Wijzigingen in SolvCRM worden automatisch verwerkt in e-Boekhouden.nl.
\end{itemize}

Daarnaast moet synchronisatie inzichtelijk zijn:
\begin{itemize}
    \item Logging van koppelingen, aanmaakacties, fouten en wijzigingen.
\end{itemize}

\subsubsection*{Factuurmutaties}
De analyse toont aan dat facturen één van de meest foutgevoelige onderdelen zijn van de huidige werkwijze. Hiervoor zijn de volgende eisen vastgesteld:

\begin{itemize}
    \item Automatische synchronisatie van verzonden facturen.
    \item Mapping van factuurmutaties.
    \item Ondersteuning voor koppelen van bestaande mutaties.
    \item Automatische aanmaak van nieuwe mutaties in beide systemen.
    \item Logging van koppelingen, aanmaakacties en fouten.
\end{itemize}

\subsubsection*{Creditnota's}
Uit gesprekken met medewerkers blijkt dat creditnota’s vaak manueel worden ingevoerd, wat foutgevoelig is. Daarom gelden de volgende eisen:

\begin{itemize}
    \item Automatische synchronisatie van creditnota’s.
    \item Mapping van creditnota’s tussen beide systemen.
    \item Mogelijkheid om bestaande creditnota’s te koppelen.
    \item Automatische aanmaak van creditnota’s.
    \item Correcte koppeling van creditnota’s aan hoofd-facturen in SolvCRM.
    \item Automatische bijwerking van factuurstatussen.
    \item Logging van koppelingen, aanmaakacties en fouten.
\end{itemize}

\subsubsection*{Betalingsmutaties}
De analyse toont dat realtime inzicht in betalingen voor klanten essentieel is. De volgende eisen zijn daarom vastgesteld:

\begin{itemize}
    \item Automatische synchronisatie van betalingsmutaties.
    \item Mapping tussen beide systemen.
    \item Mogelijkheid om bestaande betalingen te koppelen.
    \item Aanmaken van nieuwe betalingsmutaties in beide systemen.
    \item Automatische koppeling van betalingen aan de juiste factuur.
    \item Automatische bijwerking van de factuurstatus.
    \item Logging van koppelingen, aanmaakacties en fouten.
\end{itemize}

\subsection{Niet-functionele eisen}

Uit de best practices van API-integratie komen de volgende eisen:

\begin{itemize}
    \item \textbf{Betrouwbaarheid:} synchronisatie moet fouttolerant zijn.
    \item \textbf{Prestaties:} synchronisatie mag de performance van SolvCRM niet negatief beïnvloeden.
    \item \textbf{Beveiliging:} API-sleutels en gegevens moeten veilig worden opgeslagen.
    \item \textbf{Onderhoudbaarheid:} logging moet voldoende detail bieden om fouten te reproduceren.
    \item \textbf{Gebruiksvriendelijkheid:} duidelijke instellingenpagina voor synchronisatie.
\end{itemize}

\section{Resultaten voor Deelvraag 3}
% TIP: Vergelijking van integratie-opties.

\section{Resultaten voor Deelvraag 4}
% TIP: Testresultaten tijdsbesparing + foutreductie.
